\documentclass[12pt]{article}
\usepackage[utf8]{inputenc}
\usepackage[T1]{fontenc}
\usepackage[margin=1in]{geometry}
\usepackage{setspace}
\usepackage{xcolor}
\usepackage[colorlinks=true, urlcolor=blue!60!black, linkcolor=blue!60!black]{hyperref}
\setlength{\parindent}{0em}
\setlength{\parskip}{0.6em}
\renewcommand{\baselinestretch}{1.0}
\begin{document}
\begin{center}
{\LARGE \textit{Cuyler v.\ Adams}, 449 U.S.\ 433 (1981)\par}
\vspace{0.5em}
{\large Case Brief --- Two-Page Edition\par}
\end{center}

\section{Memory Jogger}
Congressional consent, even in advance, makes an interstate compact on a proper federal subject federal law; the Interstate Agreement on Detainers preserves prisoners' extradition-hearing rights (\href{https://supreme.justia.com/cases/federal/us/449/433/}{Justia}).

\section{Detailed Case Facts}
John Adams, serving 30 years in Pennsylvania, was sought by New Jersey prosecutors through a request for temporary custody under Article IV of the Interstate Agreement on Detainers (IAD). Both States had also adopted the Uniform Criminal Extradition Act, which guarantees a pre-transfer hearing. Adams was transferred without one.

\section{Procedural History}
Adams sued under \S~1983; the district court dismissed. A Pennsylvania court separately held that IAD transferees have no hearing right. The Third Circuit held the IAD is federal law and requires the hearing. The Supreme Court affirmed.

\section{Judicial Votes}
Affirmed 6--3. Brennan wrote for White, Marshall, Blackmun, Powell, and Stevens; Rehnquist dissented, joined by Burger and Stewart.

\section{Holding}
Where Congress authorizes a cooperative agreement on a subject appropriate for federal legislation, its consent makes the agreement federal law. The 1934 Crime Control Consent Act did so for the IAD, and Article IV(d) preserves the prisoner's pre-transfer hearing.

\section{Analysis of Opinions}
\textbf{Majority:} although consent is required only for agreements that threaten federal supremacy, consent actually given federalizes a compact; Article IV(d)'s ``any right'' language and its single express exception show that other extradition safeguards survive. \textbf{Rehnquist:} the ruling is ``judicial alchemy''; unnecessary consent cannot turn state law into federal law, and the decision threatens to become a ``judicial Midas.''

\section{Examples: Future Applications}
\textbf{Same-Side:} a congressionally consented disaster-response compact whose meaning is a federal question. \textbf{Opposite-Side:} an unconsented occupational-licensing compact, which stays state law. \textbf{Fence-Sitter:} whether the Uniform Act to Secure the Attendance of Witnesses is federalized by the 1934 consent.

\section{Critique}
Critics say treating a blanket advance consent as federal lawmaking sidesteps bicameralism and presentment, and the Court left unclear which reciprocal state statutes count as compacts.

\section{Key Quotations}
\begin{quote}
``But where Congress has authorized the States to enter into a cooperative agreement, and where the subject matter of that agreement is an appropriate subject for congressional legislation, the consent of Congress transforms the States' agreement into federal law under the Compact Clause.''
\end{quote}
\noindent --- Justice Brennan, for the Court, 449 U.S.\ at 440.

\end{document}
